\documentclass[12pt]{article}
\usepackage[utf8]{inputenc}
\usepackage[T1]{fontenc}
\usepackage[margin=1in]{geometry}
\usepackage{setspace}
\usepackage{xcolor}
\usepackage[colorlinks=true, urlcolor=blue!60!black, linkcolor=blue!60!black]{hyperref}
\setlength{\parindent}{0em}
\setlength{\parskip}{0.6em}
\renewcommand{\baselinestretch}{1.0}
\begin{document}
\begin{center}
{\LARGE \textit{Brown v.\ Board of Education (Brown I)}, 347 U.S.\ 483 (1954)\par}
\vspace{0.5em}
{\large Case Brief --- Two-Page Edition\par}
\end{center}

\section{Memory Jogger}
``Separate but equal'' has no place in public education: state-sanctioned school segregation violates the Equal Protection Clause (\href{https://supreme.justia.com/cases/federal/us/347/483/}{Justia}).

\section{Detailed Case Facts}
Four consolidated class actions by Black schoolchildren from Kansas, South Carolina, Virginia, and Delaware challenged laws requiring or permitting segregated schools. In Topeka, the Black and white schools were found substantially equal, so the challenge targeted segregation itself. The NAACP Legal Defense Fund, led by Thurgood Marshall, brought the cases after earlier graduate-school victories.

\section{Procedural History}
Three district courts upheld segregation under \textit{Plessy}; Delaware's courts ordered admission for tangible inequality. After argument in 1952 and reargument in 1953, the Court decided the cases on May 17, 1954, and set remedies for reargument (\textit{Brown II}).

\section{Judicial Votes}
9--0. Chief Justice Warren wrote for a unanimous Court; there were no concurrences or dissents.

\section{Holding}
Segregating public school children solely by race violates equal protection even when tangible facilities are equal. The holding reaches \textit{de jure} segregation in public education.

\section{Analysis of Opinions}
Warren found the 1868 history inconclusive and held that the Court ``cannot turn the clock back to 1868''; education must be judged by its present role in American life. Because facilities were equal, the question was whether segregation itself harms. Building on \textit{Sweatt} and \textit{McLaurin}, Warren found that it generates a feeling of inferiority, a finding Footnote 11 called ``amply supported by modern authority.'' The Court's unanimity was deliberate, to deny the South a fractured ruling to resist.

\section{Examples: Future Applications}
\textbf{Same-Side:} segregated but identical city parks or buses are invalid. \textbf{Opposite-Side:} imbalance caused only by private housing choices, without state intent, is not a violation. \textbf{Fence-Sitter:} a voluntary plan using race as a tiebreaker to integrate schools, which divided the Court in \textit{Parents Involved}.

\section{Critique}
Herbert Wechsler argued the opinion lacked a neutral principle; Charles Black answered that segregation was a caste system. Derrick Bell attributed the result to interest convergence, and originalists debate whether the 1868 meaning supports \textit{Brown}.

\section{Key Quotations}
\begin{quote}
``Today, education is perhaps the most important function of state and local governments. Compulsory school attendance laws and the great expenditures for education both demonstrate our recognition of the importance of education to our democratic society.''
\end{quote}
\noindent --- Chief Justice Warren, for a unanimous Court, 347 U.S.\ at 493.

\end{document}
